\documentclass[a4paper,11pt]{article}
\usepackage[utf8]{inputenc}
\usepackage[french]{babel}
\usepackage{amsmath}
\usepackage[osf,sups]{Baskervaldx} % lining figures
\usepackage[bigdelims,cmintegrals,vvarbb,baskervaldx]{newtxmath} % math font
\usepackage{graphicx}
\usepackage{xcolor}
\usepackage{url}
\usepackage[hidelinks]{hyperref}

% hyperrref configuration: remove ugly boxes and colors, keep urls in blue
\hypersetup{colorlinks,linkcolor=black,citecolor=black,urlcolor={blue!80!black}}


% paragraphs
\setlength{\parindent}{0em}
\setlength{\parskip}{1em}

% margins
\addtolength{\hoffset}{-4em}
\addtolength{\textwidth}{8em}
\addtolength{\voffset}{-5em}
\addtolength{\textheight}{10em}

% macros
\def\ud{\mathrm{d}}

\begin{document}
\thispagestyle{empty}

{\bf
	%Calcul numérique des fonctions de Mathieu et de leurs racines.
	%De la loi de Snell-Descartes aux symboles de Christoffel.
	Plongement des variétés de Jacobi-Maupertuis
}

{\bf Encadrants}:\\
Enric Meinhardt-Llopis \verb+<enric.meinhardt@ens-paris-saclay.fr>+\\
Léo Nicollier \verb+<leo.nicollier@ens-paris-saclay.fr>+

{\bf Objectif.}
L'objectif du stage est de construire des surfaces de révolution
en~$\mathbf{R}^3$ dont les géodésiques correspondent à les trajectoires d'un
système mécanique avec un potentiel central (orbites de Kepler, oscillations de
Hooke, etc).  Plus précisément, étant donnée une métrique radiale et conforme
sur le plan
\[
	g = g\left(\sqrt{x^2+y^2}\right)(\ud x^2+\ud y^2)
\]
il faut trouver une surface paramétrée~$\varphi:\mathbf{R}^2\to\mathbf{R}^3$
qui soit isométrique à la variété abstraite définie par~$g$.  Formellement,
cette condition revient à un système d'EDP non-linéaires de premier ordre
\[
	\varphi_x \cdot\varphi_x = g
	\qquad
	\quad
	\varphi_y \cdot\varphi_y = g
	\qquad
	\quad
	\varphi_x \cdot\varphi_y = 0
\]
que l'on devra résoudre pour trouver le plongement.  Pour les surfaces de
révolution,~$\varphi$ est déterminée par une seule
fonction~$f:\mathbf{R}\to\mathbf{R}$ et l'EDP dévient une EDO pour~$f$.
D'autre côté,
selon le principe de Maupertuis, les trajectoires d'un système physique avec
potentiel~$V(\mathbf{q})$ et énergie totale~$E$ sont, à paramétrisation près,
les géodésiques de la métrique de
Jacobi~$g=\sqrt{E-V(\mathbf{q})}\ud\mathbf{q}^2$.  Ainsi, résoudre le problème du
plongement permettra de voir les trajectoires de tout système physique avec un
potentiel radial comme les géodésiques d'une surface concrète de l'espace.
On ne comprend pas pourquoi ces surfaces n'apparaissent pas sur la couverture
de tous les livres de mécanique.
%L'objectif du stage est de prendre la loi de réfraction de Snell-Descartes
%\begin{equation}\label{eq:snell}
%\frac{\sin\theta_1}{v_1}=\frac{\sin\theta_2}{v_2}
%\end{equation}
%la manipuler, d'une manière ou d'une autre, et arriver aux équations des
%géodésiques sur une variété~\cite{vdg}
%\begin{equation}\label{eq:geo}
%\ddot q^k = \Gamma^k_{ij}\,\dot q^i\,\dot q^j
%\end{equation}
%Les équations~(\ref{eq:snell}) et~(\ref{eq:geo}) décrivent le même phénomène
%physique~: la trajectoire la plus rapide dans un milieu non homogène.  Ainsi,
%il est naturel qu'elles soient équivalentes.  En effet, l'application
%successive de la loi de Snell-Descartes sur des couches de plus en plus fines
%amène à une équation différentielle~\cite{video2}.
%%pour en avoir une visualisation voir
%Les solutions de cette équation se retrouvent être les géodésiques d'une
%variété Riemannienne dont la métrique est induite par l'indice de réfraction.
%L'objectif du stage est ainsi d'éteindre ce résultat au cas des métriques
%non-conformes (vitesse de propagation différente selon la direction) et
%non-constantes au long d'une dimension.

%\begin{tabular}{ccccc}
%	\includegraphics[width=0.20\linewidth]{f/snell.png}&
%	\includegraphics[width=0.25\linewidth]{f/snellN.png}&
%	\includegraphics[width=0.25\linewidth]{f/snellC2.png}&
%	$\color{red}\cdots$&
%	\includegraphics[width=0.18\linewidth]{f/geodesicsn.png}\\
%	$\frac{\sin\theta_1}{v_1}=\frac{\sin\theta_2}{v_2}$&
%	$\frac{\sin\theta_i}{v_i}=K$&
%	$\frac{\sin\arctan\frac{dx}{dy}}{v(y)}=K$&
%	{\color{red}\bf ?} &
%	$\ddot q^k = \Gamma^k_{ij}\,\dot q^i\,\dot q^j$
%\end{tabular}

{\bf Contexte.}
Le principe de moindre action de Maupertuis
%L'analogie entre optique et mécanique date des temps de Newton et
%Bernouilli~\cite{broer}, qui l'a utilisée pour résoudre le problème de la
%brachistocrhone (la forme de la pente qui permet d'y glisser avec le moindre
%temps).  Ainsi, la forme optimale correspond à la trajectoire d'un rayon
%lumineux sur un milieu non-homogène avec vitesse de propagation~$v=1/\sqrt{y}$.
%Le même raisonnement permet de retrouver les géodésiques sur le demi-plan de
%Poincaré~($v=1/y$) ou encore le tir parabolique~($v=y/\sqrt{1+y^2}$).  Par
%passage à coordonnées polaires, cela permettra aussi de caractériser le
%mouvement sous un potentiel central.  Si bien l'extension de la loi de Snell
%aux métriques non-isotropiques, ou même Finslériennes~\cite{markvorsen}, est un
%calcul élémentaire, son application au calcul général de géodésiques est---à
%notre connaissance---un problème ouvert.


{\bf Déroulement.}
Ce stage s'adresse aux étudiants qui souhaitent approfondir leurs connaissances
en géométrie Riemannienne et en calcul des variations.  Selon les préférences
des étudiants, on pourra faire plus ou moins des simulations par ordinateur
afin de vérifier numériquement les résultats trouvés.


\vspace{-1.5em}
\renewcommand{\refname}{\normalsize Références}
%
\begin{thebibliography}{99}
\vspace{-1em}
{\scriptsize
%\bibitem{icaza}
%	de~Icaza,~M.
%	{\it Galileo, Bernoulli, Leibniz and Newton around the brachistocrhone
%		problem},
%	Rev.~Mex.~Phys.,~1994
%\bibitem{erlichson}
%	Erlichson,~E.
%	{\it Johann Bernoulli's brachistocrhone solution using Fermat's
%		principle of least time},
%	Eur.~J.~Phys.,~1999
\bibitem{broer}
	Broer,~H.~W.
	{\it Bernoulli's light ray solution of the brachistocrhone problem
		through Hamilton's eyes},
	Int.~J.~Bifurc.~Chaos,~2014
\bibitem{video}
	3Blue1Brown
	{\it The Brachistochrone, with Steven Strogatz},
	\url{https://youtu.be/Cld0p3a43fU},~2016
\bibitem{video2}
	Veritasium
	{\it The Closest We’ve Come to a Theory of Everything},
	\url{https://youtu.be/Q10_srZ-pbs},~2024
\bibitem{vdg}
	Needham,~T.
	{\it Visual differential geometry and forms.}
	Princeton University Press,~2021
\vspace{-0.5em}
\bibitem{markvorsen}
	Markvorsen,~S., et al.
	{\it Snell's law revisited and generalized via Finsler geometry},
	Int.~J.~Geom.~Methods~Mod.~Phys.,~2023
}
\end{thebibliography}

\cleardoublepage
\Large
\[
	\color{blue}
	\frac{\sin\theta_1}{v_1}=\frac{\sin\theta_2}{v_2}
	\qquad
	%\ddot q^k = \Gamma^k_{ij}\dot q^i\dot q^j
	\qquad
	\color{red}
	\frac{\ud^2 x^k}{\ud t^2}
	=
	\Gamma^k_{ij}
	\,
	\frac{\ud x^i}{\ud t}
	\,
	\frac{\ud x^j}{\ud t}
\]
\[
	\color{blue}
	\frac{\sin\theta_1}{v_1}
	=
	\frac{\sin\theta_2}{v_2}
	=
	\frac{\sin\theta_3}{v_3}
\]
\[
	\color{blue}
	\frac{\sin\theta_1}{v_1}
	=
	\frac{\sin\theta_2}{v_2}
	=
	\frac{\sin\theta_3}{v_3}
	=
	K
\]
\[
	\color{blue}
	\frac{\sin\theta_i}{v_i}
	=
	K
\]
\[
	\color{blue}
	\frac{\sin\theta(x)}{v(x)}
	=
	K
\]
\[
	\color{blue}
	\frac{\sin\arctan\frac{dy}{dx}}{v(x)}
	=
	K
\]
\[
	\color{blue}
	\frac{y'(x)}{v(x)\sqrt{1+y'(x)^2}}
	=
	K
\]


\end{document}


% vim:set tw=79 spell spelllang=fr:
